\section*{第 \ref{ch:b2:batteries-electrolysis} 章——电池、燃料电池与电解}

\begin{solution}{exo:b2:batteries-electrolysis:1}
锌单独时约为 \qty{-0.82}{V}，电流为 0.15（模型单位）。与
铜接触时约为 \qty{-0.74}{V}，电流为 2.4。电流，进而溶解速率，
接触铜时约大十六倍。
\end{solution}

\begin{solution}{exo:b2:batteries-electrolysis:2}
$m = 500 \times 3600 \times 65.38/(2 \times 96\,485) = \qty{610}{g}$。
\end{solution}

\begin{solution}{exo:b2:batteries-electrolysis:3}
$E^\circ = \qty{2.04}{V}$（同本章）；六节电池串联，
得 \qty{12}{V}。
\end{solution}

\begin{solution}{exo:b2:batteries-electrolysis:4}
$1/6.94 = \qty{0.144}{mol}$；$0.144 \times 96\,485 = \qty{1.39e4}{C}$，即
\qty{3.86}{A.h}。
\end{solution}

\begin{solution}{exo:b2:batteries-electrolysis:5}
按 100~\%计：$200 \times 86\,400 \times 63.546/(2 \times 96\,485) = \qty{5.69}{kg}$；
$\eta_F = 5.50/5.69 = 0.967$。
\end{solution}

\begin{solution}{exo:b2:batteries-electrolysis:6}
$w = 2 \times 96\,485 \times 3.3/(0.90 \times 0.06538) = \qty{1.08e7}{J/kg}$，即
\qty{3.0}{kW.h/kg}。
\end{solution}

\begin{solution}{exo:b2:batteries-electrolysis:7}
$U = E - rI$：$r = (1.55 - 1.40)/(0.60 - 0.10) = \qty{0.30}{\ohm}$；$E = 1.55 + 0.30
\times 0.10 = \qty{1.58}{V}$。
\end{solution}

\begin{solution}{exo:b2:batteries-electrolysis:8}
阳极 \ce{2Cl- -> Cl2 + 2e-}；阴极 \ce{2H2O + 2e- -> H2 + 2OH-}。$\Delta_r G^\circ =
2(-157.24) - 2(-131.23) - 2(-237.13) = \qty{422.2}{kJ/mol}$，$U_{\min} =
\qty{2.19}{V}$。氯气会与氢氧根反应（\ce{Cl2 + 2OH- -> Cl- +
ClO- + H2O}），使两种产物都受到破坏，氯气与氢气的混合物
可能发生爆炸。
\end{solution}

\begin{solution}{exo:b2:batteries-electrolysis:9}
\ce{NaCl -> Na + 1/2Cl2}，$n = 1$：$U_{\min} = 310\,674/96\,485 = \qty{3.22}{V}$。
每千克钠：$96\,485 \times 3.22/0.02299 = \qty{1.35e7}{J}$，即
\qty{3.75}{kW.h}。
\end{solution}

\begin{solution}{exo:b2:batteries-electrolysis:10}
在锌阴极上，\ce{H+} 的还原很慢，其势垒位于
锌的波之下：电势降低时，锌先沉积，
氢气很少。在铂上，\ce{H+} 的还原很快，且
在 \qty{0}{V} 附近即开始：起初只放出氢气。铂被
锌覆盖后，阴极表现为锌阴极，锌开始沉积。
\end{solution}

\begin{solution}{exo:b2:batteries-electrolysis:11}
在 \qty{1250}{K} 时：$\Delta_f G^\circ(\ce{Al2O3}) = -1328.3 + 0.75 \times 66.0 =
\qty{-1278.8}{kJ/mol}$，$\Delta_f G^\circ(\ce{CO2}) = \qty{-396.1}{kJ/mol}$。$\Delta_r G^\circ =
3(-396.1) - 2(-1278.8) = \qty{1369}{kJ}$，$n = 12$：$U_{\min} = \qty{1.18}{V}$。惰性阳极时：
$\Delta_r G^\circ = \qty{2557.5}{kJ}$，$U_{\min} = \qty{2.21}{V}$；燃烧的碳
提供约 \qty{1.03}{V} 的功。
\end{solution}

\begin{solution}{exo:b2:batteries-electrolysis:12}
每摩尔水的电荷 $2F$ 两种情况相同：效率 $0.70/1.80 =
0.39$。其余能量转化为热量：两个氧电极的\omterm{def:b2:current-potential-curves:overpotential}{过电势}（两个方向都很慢
）、氢电极的\omterm{def:b2:current-potential-curves:overpotential}{过电势}以及欧姆损耗。
\end{solution}

\begin{solution}{pb:b2:batteries-electrolysis:1}
\textbf{1.} 铝 $+3 \to 0$；氧保持 $-2$；碳 $0 \to +4$。
\textbf{2.} 阴极 \ce{Al^3+ + 3e- -> Al}；阳极 \ce{C + 2O^2- -> CO2 + 4e-}。
\textbf{3.} \ce{2Al2O3 + 3C -> 4Al + 3CO2}，$n = 12$。
\textbf{4.} \ce{Al2O3}：$-1328.3 + 0.75 \times (1328.3 - 1262.3) = \qty{-1278.8}{kJ/mol}$；
\ce{CO2}：\qty{-396.1}{kJ/mol}。
\textbf{5.} $3(-396.1) - 2(-1278.8) = \qty{+1369}{kJ}$。
\textbf{6.} $1\,369\,000/(12 \times 96\,485) = \qty{1.18}{V}$。
\textbf{7.} $2 \times 1278.8/(12 \times 96.485) = \qty{2.21}{V}$：碳的氧化，
本身是\omterm{def:b2:reaction-free-energy:exergonic}{放能的}，承担了分解氧化物所需功的近一半。
\textbf{8.} $10^6/26.982 = \qty{3.706e4}{mol}$。
\textbf{9.} $3 \times 96\,485 \times 3.706 \times 10^4 = \qty{1.073e10}{C}$。
\textbf{10.} $1.073 \times 10^{10}/0.94 = \qty{1.141e10}{C}$。
\textbf{11.} $1.141 \times 10^{10}/3.0 \times 10^5 = \qty{3.80e4}{s}$，\qty{10.6}{h}。
\textbf{12.} $24/10.6 = 2.27$：每天 \qty{2.27}{t}。
\textbf{13.} 溶于\omterm{def:b2:batteries-electrolysis:electrolysis}{电解}质中的部分铝到达阳极气体处，并
被二氧化碳再氧化，因此其电荷被消耗了两次；另有部分电流
未经\omterm{def:b2:batteries-electrolysis:electrolysis}{电解}即穿过\omterm{def:b2:batteries-electrolysis:electrolysis}{电解}质泄漏。
\textbf{14.} $1.141 \times 10^{10} \times 4.1 = \qty{4.68e10}{J}$，\qty{13.0}{MW.h}。
\textbf{15.} \qty{13.0}{kW.h/kg}。
\textbf{16.} $3 \times 96\,485 \times 1.18/0.026982 = \qty{1.27e7}{J/kg}$，\qty{3.5}{kW.h/kg}。
\textbf{17.} 转化为热量：电极的\omterm{def:b2:current-potential-curves:overpotential}{过电势}，以及主要的
\omterm{def:b2:batteries-electrolysis:electrolysis}{电解}质欧姆压降。这些热量使\omterm{def:b2:batteries-electrolysis:electrolysis}{电解}质保持在约
\qty{1250}{K} 的熔融状态，否则将需要燃料。
\textbf{18.} $4.1 \times 3.0 \times 10^5 = \qty{1.23}{MW}$。
\textbf{19.} $28\,690/0.026982 = \qty{1.06e6}{J/kg}$，\qty{0.30}{kW.h/kg}：约占
\omterm{def:b2:batteries-electrolysis:electrolysis}{电解}能量的 2~\%。
\textbf{20.} $1.18/4.1 = 0.29$。
\textbf{21.} $3.706 \times 10^4 \times 3/4 \times 12.011 = \qty{334}{kg}$。
\textbf{22.} $3.706 \times 10^4 \times 3/4 \times 44.009 = \qty{1.22e3}{kg}$。
\textbf{23.} 阳极的高温顶部在空气中燃烧；部分碳以
一氧化碳形式离开（该温度下 \ce{C + CO2 -> 2CO}），这使
每个氧原子所耗用的碳是两倍。
\textbf{24.} 每千克铝产生 \qty{1.22}{kg} 的 \ce{CO2}。
\textbf{25.} 最小电压将升高约 \qty{1.03}{V}，但阳极
将放出氧气而非二氧化碳，且不会被消耗。
\textbf{26.} 每千克铝，该电池耗用 $\boldsymbol{\approx \qty{13}{kW.h}}$ 的
电能。
\end{solution}
